%=============================================================================!
% 1.3  ACCELERATION                                                           !
%=============================================================================!
\section{Acceleration}
\label{sec:mo-acceleration}

Because the velocity of the ball changes during its flight, a further
quantity is needed to describe how quickly it changes. This is the
acceleration, the rate of change of velocity, defined from the velocity in
exactly the way that the velocity was defined from the position,
\begin{equation*}
   a(t) = \lim_{\tau\to0}\frac{v(t+\tau) - v(t)}{\tau}
   = \frac{\dd v}{\dd t} = \frac{\dd^2 x}{\dd t^2}.
\end{equation*}
The last form records that the acceleration is obtained by differentiating
the position twice. Derivatives with respect to time occur so often that
they are also written with dots over the symbol, one dot for each
derivative: $\dot{x} = \dd x/\dd t = v$ and $\ddot{x} = \dd^2x/\dd t^2 =
a$. The dimension of acceleration is that of a velocity divided by a time,
$\mathsf{L}\mathsf{T}^{-1}/\mathsf{T} = \mathsf{L}\mathsf{T}^{-2}$, and it
is measured in metres per second per second, \si{\metre\per\second
\squared}. For the ball, differentiating $v(t) = v_0 - gt$ term by term
gives zero for the constant $v_0$ and $-g$ for the term $-gt$, so
\begin{equation*}
   a(t) = -g
\end{equation*}
at every instant of the flight. \Cref{fig:mo-ball} shows the height, the
velocity and the acceleration on a common time axis.

\begin{figure}[tb]
\centering
\includegraphics{ch01_motion/ball}
\caption{The thrown ball, $v_0 = \SI{12}{\metre\per\second}$. (a) Height,
(b) velocity and (c) acceleration against time. At the top of the flight
(dotted line) the velocity is zero and the acceleration is $-g$, as at
every other instant. The shaded area under $v(t)$ up to the top equals the
greatest height (\cref{sec:mo-integrating}).}
\label{fig:mo-ball}
\end{figure}

\subsection*{The top of the flight}

The ball reaches the top of its flight when it stops rising, that is, when
its velocity is zero. Setting $v(t) = 0$ and solving for $t$,
\begin{equation*}
   v_0 - g t = 0
   \quad\Longrightarrow\quad
   t = \frac{v_0}{g} = \frac{12}{9.80665}\,\si{\second}
   = \SI{1.224}{\second},
\end{equation*}
where the arrow $\Longrightarrow$ reads `which implies': the equation on
its right follows from the one on its left.
At this instant the ball is momentarily at rest, yet its acceleration is
still $-g$. There is no contradiction here, because a function can vanish
at an instant without its derivative vanishing at that instant. The
simplest example is $f(t) = t$, which is zero at $t = 0$ while its slope
there is 1. In the same way the velocity $v(t) = v_0 - gt$ is zero at $t =
v_0/g$, while its slope there, the acceleration, is $-g$. It is this
non-zero acceleration that carries the velocity from positive values,
through zero, to negative values, and it therefore turns the ball round.

\begin{keyidea}
Velocity is the rate of change of position, and acceleration is the rate
of change of velocity. Either can be zero at an instant at which the other
is not.
\end{keyidea}

\subsection*{Speeding up and slowing down}

The signs of $v$ and $a$ together determine whether a motion is speeding
up or slowing down. The speed is the size of the velocity, $\lvert
v\rvert$, and its square is $\lvert v\rvert^2 = v^2$. Since the speed is
never negative, it grows exactly when its square grows, and the square
spares us the absolute value. To see the rule, we therefore ask how $v^2$
changes in time. The square is a function of $v$, which is itself a
function of $t$, and the chain rule (below) gives
\begin{equation*}
   \frac{\dd}{\dd t}\,v^2 = 2v\,\frac{\dd v}{\dd t} = 2va .
\end{equation*}
When $v$ and $a$ have the same sign their product is positive, $v^2$
grows and the speed increases; when their signs differ the product is
negative and the speed decreases. While the ball rises, $v > 0$ and $a <
0$, so it slows down. Once it begins to fall, $v$ and $a$ are both
negative, and it speeds up.

\begin{toolbox}[The chain rule]
Let $f$ depend on $t$ through an intermediate quantity $u(t)$. Over a short
interval $\tau$, $u$ changes by $\Delta u = u(t+\tau) - u(t)$, and
multiplying and dividing by $\Delta u$,
\begin{equation*}
   \frac{f(u + \Delta u) - f(u)}{\tau}
   = \frac{f(u + \Delta u) - f(u)}{\Delta u}\times\frac{\Delta u}{\tau} .
\end{equation*}
As $\tau \to 0$, $\Delta u \to 0$ too, so the first factor approaches
$f'(u)$, the derivative of $f$ with respect to $u$, and the second
approaches $\dd u/\dd t$. At $\Delta u = 0$, define the first factor by
its limiting value $f'(u)$; its product with $\Delta u/\tau$ is then zero,
as is the left side. This also covers intervals on which $u$ does not
change. Hence
\begin{equation*}
   \frac{\dd}{\dd t}\,f\bigl(u(t)\bigr) = f'(u)\,\frac{\dd u}{\dd t}.
\end{equation*}
With $f(u) = u^2$, so that $f'(u) = 2u$, and $u = v$, this gives the rate
$2v\,\dd v/\dd t$ used above.
\end{toolbox}

\notebook{01}{move the time cursor along $x$, $v$ and $a$}
